\documentclass[11pt]{article}
\usepackage[utf8]{inputenc}
\usepackage{amsmath,amssymb,amsthm}
\usepackage[margin=1in]{geometry}
\usepackage{hyperref}
\usepackage{xcolor}
\usepackage{enumitem}
\usepackage{newunicodechar}
\newunicodechar{ℝ}{\ensuremath{\mathbb{R}}}

\newcommand{\userbox}[1]{\par\medskip\begin{quote}\small\textbf{\textcolor{blue!60}{DM:}} #1\end{quote}\medskip}
\newcommand{\claudebox}[1]{\par\medskip\begin{quote}\small\textbf{\textcolor{orange!60!black}{Claude:}} #1\end{quote}\medskip}
\newcommand{\gptbox}[1]{\par\medskip\begin{quote}\small\textbf{\textcolor{green!50!black}{GPT-5.5 Pro:}} #1\end{quote}\medskip}

\title{From Primer to Lean:\\A retrospective on formalising the\\anharmonic Laplace asymptotic}
\author{Daniel Murfet}
\date{April 2026}

\begin{document}
\maketitle

\begin{abstract}
This document is a retrospective on a single morning of formalisation work
in Lean 4 + Mathlib, in which the SLT Susceptibility Primer's anharmonic
covariance formula
\[
  \operatorname{Cov}_t[x^2,x]
  \;=\; -\frac{2\alpha}{\lambda^3 t^2} + o(t^{-2}) \qquad (t\to\infty)
\]
was formalised end-to-end, producing 3279 lines of Lean across ten files,
60+ proved theorems, and zero \texttt{sorry}, \texttt{axiom},
or \texttt{native\_decide} occurrences. The collaboration involved Claude
(Anthropic, Opus 4.7) as the primary working partner with three targeted
consultations to GPT-5.5 Pro at strategic decision points. The formalisation
borrowed structure and tooling from Geoffrey Irving's
\texttt{aks} repository. End-to-end the session took 5h 44m of wall-clock
time (2h 30m of API time) and \$198.53 of model usage. The aim of this
retrospective is to record honestly what proportion of the work was Mathlib
lookup, what was new, where progress stalled, and how the multi-model
workflow contributed.
\end{abstract}

\tableofcontents

\section{Setting the stage}

In a working note dated 27 April 2026 (``Formalisation First Steps''), DM
sketched a path forward for formalisation at Timaeus, building on
conversations with Geoffrey Irving and Edmund Lau:

\begin{quote}
It's natural to go after big theorems. Geoffrey's \texttt{aks} repo is
super cool. It would be tempting for us to follow suit and go after one of
the main theorems of SLT. However that seems to have two problems: it's
unclear if the timing is right, given model capability increases and the
gaps in Mathlib; and even if we were to formalise say the free energy
formula, it's not immediately clear what that buys us in terms of forward
momentum across our research.

[\ldots] If I reflect on where the bottleneck is on using Claude/GPT to go
faster on theoretical research, it is the verification step. [\ldots] So
in this world I just tell Claude to go ``do the lower-order cases of this
kind of asymptotic expansion with Wick's theorem'' and it goes and writes
Section 4 of the primer above, verifies Lemma 20 and then just presents me
with the LaTeX and some Lean files. I then read the Lean file for Lemma 20
and just have to verify that the statement is indeed what I think and what
the LaTeX says, which takes me 20 minutes.
\end{quote}

This retrospective is the first concrete experiment along that path. The
``Lemma 20'' referred to is the one-dimensional anharmonic case of the
Laplace expansion in Section 4 of the SLT Susceptibility Primer
(Elliott--Murfet, 2026), and ``some Lean files'' is the
\texttt{timaeus-research/laplace} repository whose construction we record
here.

The motivation for picking this particular target was:
\begin{itemize}
\item The mathematics is genuinely useful but routine. It is the kind of
``boring calculation that I am 100\% sure I can do if I just sit down for
an afternoon'' but for which DM does not have a free afternoon. Formalising
it is pure win: the calculation is fun but DM is not learning anything new
by doing it himself.
\item It is well within the joint scope of Opus 4.7 + GPT-5.5 Pro, but
non-trivially so: there are sub-goals (Mill's-ratio tail bounds, Taylor
remainders, asymptotic algebra) that exercise different parts of Mathlib.
\item It can be inspected independently: anyone reading the primer can read
\texttt{cov\_anharmonic\_asymptotic} in Lean and check, in 20 minutes,
that the statement matches.
\item Building it forces the development of a small library of pieces
(scalar Taylor bounds, 1D Gaussian moments, Mill's-ratio bounds, rescaling
identities) that will be reused by subsequent work in Timaeus's research
agenda.
\end{itemize}

\section{The thing to be formalised}

The headline result of the formalisation is:

\begin{quote}\itshape
For the anharmonic potential
\[
  L(x) \;=\; \frac{\lambda}{2}x^2 + \frac{\alpha}{6}x^3 + \frac{\gamma}{24}x^4
\]
with $\lambda > 0$, $\gamma > 0$, and the discriminant condition
$\alpha^2 < 3\lambda\gamma$,
\[
  \operatorname{Cov}_t[x^2,x]
  \;=\; -\frac{2\alpha}{\lambda^3\,t^2} + o(t^{-2})
  \quad\text{as } t\to\infty,
\]
where $\operatorname{Cov}_t[\phi,\psi]$ is the covariance under the Gibbs
measure $e^{-tL(x)}\,dx / Z(t)$.
\end{quote}

In Lean (writing \texttt{R} for the type of real numbers, which the
Lean source spells $\mathbb{R}$):
\begin{verbatim}
theorem cov_anharmonic_asymptotic
    {lam alpha gamma : R}
    (hlam : 0 < lam) (hgamma : 0 < gamma) (hdisc : alpha ^ 2 < 3 * lam * gamma) :
    Filter.Tendsto
      (fun t : R =>
        t ^ 2 *
          Laplace.gibbsCov (anharmonicPotential lam alpha gamma) t
            (fun x => x ^ 2) (fun x => x))
      Filter.atTop
      (nhds (-2 * alpha / lam ^ 3))
\end{verbatim}

\paragraph{Where this lives in the primer.} This is equation (4.10) in the
primer, and it is the lowest-dimensional concrete payoff of the asymptotic
machinery there. The primer derives it via Wick contractions on the
rescaled coordinate $u = x\sqrt{\lambda t}$, expanding $e^{-s_t(u)}$ to
first order in the cubic and quartic perturbations.

\paragraph{Why the discriminant condition.} The primer's calculation is
formal: it expands $e^{-tL(x)}$ as a Gaussian times a perturbation
series. For the resulting moment integrals to be well-defined and dominated
by the quadratic minimum at $x = 0$, $L$ must be coercive (bounded below
by a positive multiple of $x^2$) globally. This is automatic in the
multivariate primer setting under generic-position assumptions, but in 1D
with a cubic term it requires the discriminant condition
$\alpha^2 < 3\lambda\gamma$. Without it, $L$ can have lower minima away
from $0$, and the Laplace asymptotic is not driven by $x = 0$.

The hypothesis was added at GPT-5.5 Pro's insistence during the strategy
consultation; see \S\ref{sec:strategy}.

\section{Background from the primer}

The primer's Section 4 computes asymptotic expansions of expectations
under $p_t(x) \propto e^{-tL(x)}$. The mechanism, in one dimension, is:

\begin{enumerate}
\item Rescale $u = x\sqrt{\lambda t}$. Then $p_t$ becomes
$\propto e^{-u^2/2 - s_t(u)}\,du$ where
\[
  s_t(u) \;=\; \frac{\alpha}{6\lambda^{3/2}}\frac{u^3}{\sqrt t}
            + \frac{\gamma}{24\lambda^2}\frac{u^4}{t}
\]
collects the cubic and quartic perturbations.
\item Expand $e^{-s_t} = 1 - s_t + \tfrac12 s_t^2 - \cdots$. To get
$O(t^{-2})$ accuracy on $\operatorname{Cov}_t[x^2,x]$ it suffices to keep
the first-order term and bound the rest.
\item Compute the leading rescaled moments
$M_k = \int u^k e^{-u^2/2}\,du = (k-1)!! \sqrt{2\pi}$ for even $k$, zero
for odd $k$.
\item Assemble: $\langle x\rangle_t, \langle x^2\rangle_t,
\langle x^3\rangle_t$ from the $J_n(t) = \int u^n e^{-u^2/2 - s_t(u)}\,du$,
the $I_n(t) = \int x^n e^{-tL(x)}\,dx$ via the substitution
$(\sqrt{\lambda t})^{n+1} I_n = J_n$, and finally
$\operatorname{Cov}_t[x^2,x] = \langle x^3\rangle_t -
 \langle x^2\rangle_t \langle x\rangle_t$.
\item The leading $t^{-2}$ coefficient is the algebraic identity
$-5\alpha/(2\lambda^3) - (1/\lambda)\cdot(-\alpha/(2\lambda^2)) =
 -2\alpha/\lambda^3$, which is real-valued ring algebra.
\end{enumerate}

This is what gets formalised. The primer treats it informally; the Lean
proof has to make every domination, every tail bound, every change of
variable, and every algebraic step explicit.

\section{Borrowings from \texttt{aks}}\label{sec:aks}

Geoffrey Irving's \href{https://github.com/girving/aks}{\texttt{aks} repository}
provided the model and several concrete artefacts:

\begin{itemize}
\item \textbf{The discipline.} \texttt{aks} formalises the
Agrawal--Kayal--Saxena primality theorem with a strict ``no
\texttt{sorry}, no \texttt{axiom}, no \texttt{native\_decide}'' rule. We
adopted the same rule for \texttt{laplace}, including the
audit script.
\item \textbf{\texttt{scripts/sorries}.} The audit tool that scans the
codebase for \texttt{sorry}, \texttt{\#exit}, \texttt{native\_decide}, and
\texttt{axiom} occurrences was lifted nearly verbatim from
\texttt{aks/scripts/sorries}. It was run before every commit.
\item \textbf{\texttt{scripts/lean-search}.} The Python wrapper around
\texttt{leansearch.net} for semantic Mathlib search was also lifted from
\texttt{aks}. It was used heavily during stage 2 (Gaussian moments) and
stage 3 (tail bounds).
\item \textbf{\texttt{CLAUDE.md} structure.} The project guidance file
follows the \texttt{aks} pattern: a short statement of the goal, a list of
build commands, conventions (toolchain pinning, no
\texttt{native\_decide}, etc.), a Mathlib API reference that grows as the
project goes, and a tactic-gotchas section. The intent --- captured by
Geoffrey's repo and adopted here --- is that an AI assistant who has not
seen the project before should be able to pick up the playbook from
\texttt{CLAUDE.md} alone.
\item \textbf{Lake configuration template.} The \texttt{lakefile.toml} and
\texttt{lean-toolchain} pinning style (matching Mathlib at the
\texttt{v4.29.0} tag) was lifted directly from \texttt{aks}.
\item \textbf{Workflow philosophy.} The aks-style discipline of
``skeleton correctness $>$ filling sorries'': a \texttt{sorry} with a
correct statement is valuable, while a \texttt{sorry} with a wrong
statement is actively misleading. We followed this throughout: at any
intermediate point the file built, with placeholder \texttt{sorry}s that
were filled in only once we were confident the statements were right.
\end{itemize}

These borrowings turned what would have been a multi-day setup (toolchain
configuration, build infrastructure, audit tooling, prompt scaffolding)
into a 30-minute setup and let the project start working on actual
mathematics from the first turn. Quoting DM's earlier note: ``By grabbing
some of Geoffrey's aks pieces, Claude + Lean is already making some
progress in 30min this morning.''

\section{Stages of the formalisation}

The work proceeded in eight rough stages. For each I record the headline
theorem(s), the proportion of work that was Mathlib lookup vs.\ new, and
where applicable the points at which DM intervened.

\subsection{Stage 0: Strategy consultation}\label{sec:strategy}

Before writing any Lean, Claude proposed three approaches to the primer's
formula:

\begin{enumerate}
\item Rescaled Gaussian + perturbation expansion (the primer's own route).
\item Watson--Erd\'elyi asymptotic expansion of parameter-dependent
integrals.
\item Morse normal form via local inverse function theorem.
\end{enumerate}

GPT-5.5 Pro was consulted on which to pick.

\gptbox{Use (1), but don't localise first. With the right hypothesis, you
can prove a global first-order perturbative bound and skip most of the
ugly tail machinery for this stage. \emph{First: the statement needs a
hypothesis fix.} As written, with only $\lambda > 0$, $\gamma > 0$,
$\alpha \in \mathbb R$, the claim is false in general: if $\alpha$ is
large enough, $L$ has lower minima away from $0$, so Laplace is not
dominated by $x = 0$. A sufficient condition for $0$ to be the unique
global minimum is $\alpha^2 < 3\lambda\gamma$.}

This was the first decisive intervention. The hypothesis fix
($\alpha^2 < 3\lambda\gamma$) made the rest of the proof tractable, and
the recommendation to do a global first-order bound (rather than a
localisation theorem in the Watson style) avoided multi-week
infrastructure work on inverse function theorems and Jacobian expansions.
GPT-5.5 Pro's strategy memo is preserved at
\texttt{gpt\_responses/strategy\_consultation.md}.

\subsection{Stage 1: Gibbs definitions and the scalar Taylor cornerstone}

\textbf{Files:} \texttt{Laplace/Gibbs.lean} (65 lines),
\texttt{Laplace/ScalarBound.lean} (179 lines).

The Gibbs definitions (\texttt{partitionFunction}, \texttt{gibbsExpectation},
\texttt{gibbsCov}) were straightforward Mathlib applications; nothing
new.

The scalar Taylor cornerstone is
\[
  \big|\, e^{-z} - (1 - z)\, \big| \;\le\; \tfrac{z^2}{2}\cdot\max(1, e^{-z}),
\]
proved as \texttt{abs\_exp\_neg\_sub\_one\_add\_le}. This is the load-bearing
inequality for the entire perturbative bound: every later $O(1/t)$
estimate ultimately reduces to multiplying both sides of this by a
Gaussian, integrating, and bounding by Gaussian moments.

\textbf{Mathlib gap:} this exact form was not in Mathlib. The closely
related \texttt{Real.add\_one\_le\_exp} ($1 + z \le e^z$) is, but it gives
only the one-sided bound. The cornerstone splits on the sign of $z$:
above zero, $e^{-z} - 1 + z \le z^2/2$ by Taylor; below zero,
$e^{-z} - 1 + z \le (z^2/2)e^{-z}$ by the same Taylor expansion of $e^{-z}$
around $z = 0$ but with the second-derivative bound shifted. Both halves
were written from scratch.

\subsection{Stage 2: 1D Gaussian moments}

\textbf{File:} \texttt{Laplace/OneD/GaussianMoments.lean} (231 lines).

The standard moments $M_{2k} = (2k-1)!!\sqrt{2\pi}$ and $M_{2k+1} = 0$.

\textbf{Mathlib provided:}
\begin{itemize}
\item \texttt{integral\_gaussian} (the base case
$\int e^{-bx^2}\,dx = \sqrt{\pi/b}$).
\item \texttt{integral\_rpow\_mul\_exp\_neg\_mul\_rpow} (general
Gamma-function form for half-line integrals).
\item \texttt{Real.Gamma\_nat\_add\_half}
($\Gamma(k+1/2) = (2k-1)!!\sqrt\pi/2^k$).
\item \texttt{Nat.doubleFactorial} and its recurrences.
\item \texttt{integral\_comp\_abs} (folding the real line to $(0,\infty)$
via $|x|$).
\end{itemize}

\textbf{What was new:} stitching these into the four user-facing forms
(\texttt{integral\_pow\_mul\_exp\_neg\_sq\_half\_Ioi},
\texttt{integral\_pow\_mul\_exp\_neg\_sq\_half},
\texttt{gaussian\_moment\_normalized}, and
\texttt{integral\_pow\_mul\_exp\_neg\_t\_sq\_half}) involved a surprising
amount of bridge work: converting between $x^{2k}$ (natural-number power)
and $x^{2k\cdot 1.0}$ (rpow), folding even powers via $|x|^{2k} = x^{2k}$,
and feeding $1/2$ into exponents without having it eaten by an over-eager
\texttt{simp}. The recurring tactic gotchas from this stage are recorded
in \texttt{CLAUDE.md} (e.g.\ \emph{cascading \texttt{rw [(1/2) = 2$^{-1}$]}
stomps inside exponents}).

\textbf{Proportion:} mathematically, all the content is Mathlib's. The
work was bookkeeping.

\subsection{Stage 3: Anharmonic potential and coercivity}

\textbf{File:} \texttt{Laplace/OneD/Anharmonic.lean} (140 lines).

The headline theorem is \texttt{anharmonic\_coercive}: under
$\alpha^2 < 3\lambda\gamma$, there exists $c > 0$ with $L(x) \ge c\,x^2$
for all $x$. Proof: complete the square inside
$\frac{\lambda}{2} + \frac{\alpha}{6}x + \frac{\gamma}{24}x^2$ and use
the discriminant.

\textbf{Mathlib provided:} \texttt{discrim} and the basic theory of
quadratic forms.

\textbf{What was new:} the application to this specific anharmonic
potential, written by Claude. About 90\% of this stage was new; only the
discriminant lemma is from Mathlib.

\subsection{Stage 4: Tail bounds and localisation}

\textbf{Files:} \texttt{Laplace/OneD/TailBound.lean} (532 lines),
\texttt{Laplace/OneD/Localisation.lean} (164 lines).

A small Mill's-ratio family: bounds of the form
\[
  \int_M^\infty x^k e^{-bx^2/2}\,dx \;\le\; (\text{polynomial in } M)
                                     \cdot e^{-bM^2/2}
\]
for $k = 0, 1, 2$ and rescaled exponents.

\textbf{Mathlib provided:} the basic Gaussian integrals and the
fundamental theorem of calculus over half-lines. The tail bounds
themselves were not in Mathlib at the form needed.

\textbf{What was new:} the entire 532-line tail-bound development. In
hindsight the localisation file was over-engineered for the eventual
proof: GPT-5.5 Pro's strategy of doing a \emph{global} first-order
remainder bound (\S\ref{sec:strategy}) made these tail bounds essentially
unused. They are kept because they will be useful for future work on
multivariate Laplace asymptotics, where global bounds are not available.

This was the first place where Claude noticeably overshot the strict
needs of the target. The second was \S\ref{sec:harmonic}.

\subsection{Stage 5: Harmonic Gibbs (an over-engineered detour)}\label{sec:harmonic}

\textbf{File:} \texttt{Laplace/OneD/Harmonic.lean} (165 lines).

Closed-form expectations under the harmonic Gibbs measure $e^{-tL_\lambda}$
with $L_\lambda(x) = \tfrac\lambda2 x^2$:
$\langle x^{2k}\rangle = \tfrac{(2k-1)!!}{(\lambda t)^k}$, odd moments
zero, and the trivial covariance $\operatorname{Cov}_t[x^2, x] = 0$ in the
harmonic case.

\textbf{Why this was a detour.} An early plan (before GPT-5.5 Pro's
strategy memo) was to compute the anharmonic case as a perturbation
\emph{around} the harmonic case in original coordinates. After the
strategy memo, the actual proof works in rescaled coordinates and never
uses the harmonic-Gibbs closed forms directly. This file is kept as a
sanity check (the rescaled-coordinate moments must agree with the
original-coordinate moments at $\alpha = \gamma = 0$) and as a piece of
infrastructure for future work.

\textbf{Mathlib provided:} all the moment integrals from stage 2.
\textbf{What was new:} the lifting to the Gibbs setting, mostly algebraic
massaging.

\subsection{Stage 6: The rescaling identity and pointwise perturbation bound}

\textbf{File:} \texttt{Laplace/OneD/Rescaling.lean} (230 lines).

The two analytic cornerstones:
\begin{enumerate}
\item \texttt{anharmonic\_rescaling\_identity}:
$tL(u/\sqrt{\lambda t}) = u^2/2 + s_t(u)$.
\item \texttt{rescaled\_max\_decay}:
$e^{-u^2/2}\cdot\max(1, e^{-s_t(u)}) \le e^{-c_0 u^2}$
under coercivity, with $c_0 = \min(1/2, c/\lambda)$.
\end{enumerate}

\textbf{Mathlib provided:} \texttt{Real.exp\_le\_exp} and basic estimates.
\textbf{What was new:} the algebraic identity (a one-line calculation
written carefully), and the decay bound (which combines the coercivity
estimate from stage 3 with the rescaling identity).

\subsection{Stage 7: The integral remainder and $J_n$ asymptotics}

\textbf{File:} \texttt{Laplace/OneD/IntegralRemainder.lean} (1541 lines,
about half the project).

This is the analytic core. The chain of theorems is:
\begin{itemize}
\item \texttt{perturbation\_remainder\_pointwise} +
\texttt{perturbation\_remainder\_combined}: pointwise
$|u^n e^{-u^2/2}(e^{-s_t} - (1 - s_t))| \le C/t \cdot |u|^n(u^6 + u^8)e^{-c_0 u^2}$.
\item \texttt{perturbation\_remainder\_integral\_bound}: integrating the
pointwise bound, $|\int \cdots| \le K/t$.
\item \texttt{linearised\_integral\_decomposition}:
$\int f(u)(1 - s_t(u))\,du = M_n - (A/\sqrt t)M_{n+3} - (B/t)M_{n+4}$.
\item \texttt{J\_n\_asymptotic}: the master $J_n$ asymptotic for general
$n$.
\item \texttt{J\_0\_asymptotic}, \texttt{J\_1\_asymptotic},
\texttt{J\_2\_asymptotic}, \texttt{J\_3\_asymptotic}: parity-specialised
specifications.
\item \texttt{I\_n\_J\_n\_relation}:
$(\sqrt{\lambda t})^{n+1} I_n = J_n$ (substitution identity).
\item \texttt{I\_0\_asymptotic}, \texttt{I\_1\_asymptotic},
\texttt{I\_2\_asymptotic}, \texttt{I\_3\_asymptotic}: original-coordinate
moment asymptotics.
\item \texttt{cov\_coefficient\_identity}:
$-5\alpha/(2\lambda^3) - (1/\lambda)\cdot(-\alpha/(2\lambda^2)) = -2\alpha/\lambda^3$.
\item Four \texttt{tendsto\_*} theorems converting the asymptotic bounds
to \texttt{Filter.Tendsto} statements.
\item \texttt{cov\_anharmonic\_J\_form\_asymptotic}: the headline result
in the rescaled $J_n$ form.
\item \texttt{cov\_anharmonic\_asymptotic}: the headline result in the
original \texttt{gibbsCov} form.
\end{itemize}

\textbf{Mathlib provided:} integration linearity (\texttt{integral\_add},
\texttt{integral\_sub}, \texttt{integral\_const\_mul}), change of
variables (\texttt{Measure.integral\_comp\_mul\_right}), the
\texttt{Filter.Tendsto} algebra and asymptotic notation
(\texttt{Asymptotics.IsBigO}, \texttt{Asymptotics.IsLittleO}).

\textbf{What was new:} essentially all of it. About 1500 lines of Lean
specific to this proof.

\subsection{The intervention that broke the J-form / gibbsCov bridge}\label{sec:bridge}

The last step --- bridging
\texttt{cov\_anharmonic\_J\_form\_asymptotic} (rescaled-$J_n$ form) to
\texttt{cov\_anharmonic\_asymptotic} (original \texttt{gibbsCov} form)
--- got stuck for several hours. The goal, after unfolding definitions
and substituting $I_n = J_n / (\sqrt{\lambda t})^{n+1}$, was a single
rational identity:
\begin{multline*}
  t^2 \cdot \frac{Z\cdot I_3 - I_2\cdot I_1}{Z^2}
  \;=\; \frac{\sqrt t \cdot (J_0 J_3 - J_2 J_1)}{\lambda\sqrt\lambda \cdot J_0^2}.
\end{multline*}

This is true. The variables are $\lambda, t, J_0, J_1, J_2, J_3$ and the
square roots $\sqrt\lambda, \sqrt t, \sqrt{\lambda t}$ which appear in
the substitution identity. Standard tactics --- \texttt{ring},
\texttt{field\_simp; ring}, \texttt{linear\_combination},
\texttt{polyrith}, \texttt{nlinarith} with the relevant
\texttt{Real.sq\_sqrt} facts --- all failed: \texttt{ring} cannot reason
about $\sqrt{\,\cdot\,}$, and the others did not find the right
combination.

After three failed attempts DM noticed:

\userbox{Seems like you might be struggling a bit, remember you can always
share your troubles with 5.5Pro.}

GPT-5.5 Pro's response was the third decisive intervention:

\gptbox{Use fresh variables for the square roots, then rewrite \texttt{lam}
and \texttt{t} as squares before \texttt{field\_simp}.
\texttt{linear\_combination}/\texttt{polyrith} are the wrong tools here;
this is a single rational identity, not a linear combination of
hypotheses. After \texttt{rw [$\leftarrow$\,hsl2, $\leftarrow$\,hst2]},
there are no sqrt axioms left: everything is in the polynomial variables
\texttt{sl, st}. Then \texttt{field\_simp} turns the goal into a genuine
polynomial identity, and \texttt{ring} closes it.}

The recipe: introduce \texttt{set sl := Real.sqrt lam},
\texttt{set st := Real.sqrt t}; prove $\texttt{sl}^2 = \lambda$,
$\texttt{st}^2 = t$; rewrite $\lambda$ and $t$ as $\texttt{sl}^2,
\texttt{st}^2$ in the goal; then \texttt{field\_simp} alone (no
\texttt{ring} needed). This worked immediately. The fix is preserved
verbatim at the end of \texttt{IntegralRemainder.lean} (lines
1527--1539).

GPT-5.5 Pro's bridge memo is preserved at
\texttt{gpt\_responses/bridge\_help.md}.

\section{Where I got stuck and how 5.5 Pro helped}

Three GPT-5.5 Pro consultations bracket the project. Each was a strategic
inflection point.

\paragraph{Strategy (\S\ref{sec:strategy}).} Picked the right approach
(rescaled Gaussian, global first-order bound) before Claude could spend
days building Watson--Erd\'elyi infrastructure. Identified the missing
hypothesis ($\alpha^2 < 3\lambda\gamma$) without which the theorem is
false. Proposed the exact form of the load-bearing scalar Taylor lemma
($e^{-z} - (1-z)$ bound) that became the cornerstone in stage 1.
\textbf{Without this consultation, the project might never have
finished.}

\paragraph{Progress check.} After completing the analytic core, Claude
proposed a generic parity framework for $J_n$ that would have added
significant bookkeeping. GPT-5.5 Pro's response:

\gptbox{Skip generic step 10/11. Specialise immediately to $n = 0,1,2,3$
and work with the rescaled unnormalized moments $J_n(t)$. This avoids a
generic parity framework and most bookkeeping. Avoid \texttt{Real.rpow}
if possible. For $n = 0,1,2,3$, write scaling in terms of
\texttt{Real.sqrt} and integer powers.}

This redirected the work away from a generic infrastructure and toward
four targeted specifications. The advice on avoiding \texttt{Real.rpow}
turned out to be especially load-bearing: \texttt{Real.rpow} interacts
badly with \texttt{nlinarith} and with the \texttt{Real.sqrt} idioms used
elsewhere in the project, and the bridge step in \S\ref{sec:bridge} would
have been substantially harder if the moment scaling were stated via
\texttt{rpow}. The full memo is at
\texttt{gpt\_responses/progress\_check.md}.

\paragraph{Bridge help (\S\ref{sec:bridge}).} The final-step unblock:
the recipe of substituting fresh symbols for the square roots and
rewriting $\lambda, t$ as squares before \texttt{field\_simp}. Without
this advice the project would have had to either (a) keep trying tactics
in the dark, (b) abandon the bridge and present only the rescaled-$J_n$
form, or (c) abandon $\lambda$-as-a-parameter and work over a fixed
$\lambda = 1$ baseline.

\paragraph{Pattern.} GPT-5.5 Pro's contributions were uniformly
\emph{strategic}: choosing the right approach, identifying the right
hypothesis, picking the right tactic. None of the consultations involved
GPT writing Lean code. Claude wrote all the Lean. GPT chose the road.

\section{What was Mathlib, what was new}

A rough accounting, in lines of Lean:

\begin{center}
\begin{tabular}{lrrl}
\hline
File & Lines & Mathlib & New content \\
\hline
\texttt{Gibbs.lean} & 65 & 100\% & Definitions only \\
\texttt{ScalarBound.lean} & 179 & 30\% & The Taylor cornerstone \\
\texttt{GaussianMoments.lean} & 231 & 80\% & Bridge work \\
\texttt{Anharmonic.lean} & 140 & 10\% & Coercivity \\
\texttt{TailBound.lean} & 532 & 5\% & Mill's-ratio family \\
\texttt{Localisation.lean} & 164 & 5\% & Tail localisation \\
\texttt{Harmonic.lean} & 165 & 60\% & Closed forms \\
\texttt{Rescaling.lean} & 230 & 5\% & Identity + decay \\
\texttt{IntegralRemainder.lean} & 1541 & 10\% & Analytic core \\
\hline
Total & 3247 & $\sim$20\% & $\sim$80\% \\
\hline
\end{tabular}
\end{center}

(``Mathlib'' here is shorthand for ``content essentially supplied by
Mathlib lemmas, modulo minor packaging''. ``New'' means proofs we wrote.
The percentages are eyeball estimates.)

\paragraph{The Mathlib gaps.} The largest gaps were:
\begin{itemize}
\item No two-sided $|e^{-z} - (1-z)|$ bound. Mathlib has
\texttt{Real.add\_one\_le\_exp} but not the symmetric version with
$\max(1, e^{-z})$.
\item No Mill's-ratio family at the strength needed
(\texttt{TailBound.lean}). Mathlib has \texttt{ProbabilityTheory.gaussian\_*}
results but not the polynomial-times-Gaussian half-line bounds.
\item No global first-order Laplace remainder lemma. The 200-line
\texttt{perturbation\_remainder\_integral\_bound} is the workhorse and
has no Mathlib analogue.
\item No off-the-shelf substitution identity
$(\sqrt{\lambda t})^{n+1} I_n = J_n$. The change-of-variables result
\texttt{Measure.integral\_comp\_mul\_right} is in Mathlib, but the
specific application here needed careful packaging.
\end{itemize}

\paragraph{Where Mathlib carried us.} \texttt{Filter.Tendsto},
\texttt{Asymptotics.IsBigO}, the integral linearity lemmas, the basic
Gaussian integral, the double-factorial / Gamma function bridge, and the
real-analysis tactic infrastructure (\texttt{positivity}, \texttt{linarith},
\texttt{nlinarith}, \texttt{ring}, \texttt{field\_simp}) were all
indispensable. About 70\% of the project would have been infeasible
without these.

\section{Where DM intervened}

Across the (single morning) session, DM's interventions fell into a
small number of categories.

\paragraph{Setting the goal.}
\begin{quote}
\itshape
\userbox{I want to try with Lean formalisation of some of the Wick
contraction / asymptotic expansions of $\exp(-tL)$ in the primer, what
do you think?}

\userbox{Where is that in the primer?}

\userbox{Ok let's get the exact formula in a Lean theorem please.}
\end{quote}

DM steered Claude away from generic asymptotic-expansion infrastructure
toward the specific primer formula, and refused to accept a partial
result that stopped short of the named covariance.

\paragraph{Sanity-checking the trajectory.}
\begin{quote}
\itshape
\userbox{Can you explain how what we formalised actually proves this
formula about $\operatorname{Cov}_t[x^2,x]$?}
\end{quote}

This came after stage 7's $I_n$ asymptotics were complete. Claude wrote
out the algebraic manipulation $\langle x^3\rangle - \langle x^2\rangle
\langle x\rangle$, which surfaced (a) the fact that the Gibbs ratios
$I_n/I_0$ had not yet been proved to admit asymptotics in the right form
and (b) the missing coefficient identity. Both were filled in.

\paragraph{Pushing through to the end.}
\begin{quote}
\itshape
\userbox{This is great. Feel free to consult with 5.5 Pro about next
steps, but I'd like us to continue.}

\userbox{Continue, let's just go all the way now.}
\end{quote}

A recurring pattern was Claude proposing to ``stop at a clean
intermediate result'' (e.g.\ stopping at the $I_n$ asymptotics; stopping
at the rescaled-$J_n$-form theorem). DM consistently declined these
exits and insisted on going all the way to a Lean theorem stated in
\texttt{gibbsCov} form.

\paragraph{The 5.5 Pro intervention.}
\begin{quote}
\itshape
\userbox{Seems like you might be struggling a bit, remember you can
always share your troubles with 5.5 Pro.}
\end{quote}

This is the prompt that produced the bridge fix in \S\ref{sec:bridge}.
DM noticed Claude was thrashing on the same goal across multiple
attempts (the signs from \texttt{CLAUDE.md}: ``oscillating between
approaches, growing helper count without progress, repeated
restructuring''), and suggested escalating to GPT-5.5 Pro before sinking
more time. It is the kind of intervention that an experienced research
supervisor learns to make for graduate students: stop, look up, ask for
help.

\paragraph{Categories of intervention.} Mapping the pattern from the
\texttt{coevolving\_coefficients} retrospective onto this project:

\begin{itemize}
\item \emph{Directive} (telling Claude exactly what to write): rare here.
The directive moments were almost all setup-related (``borrow
Geoffrey's audit script'', ``pin Mathlib at v4.29.0'', ``add a
\texttt{CLAUDE.md}'').
\item \emph{Exploratory} (open-ended questions): also rare. The most
notable was ``How much of this is background material that was
previously missing in mathlib, versus tools that are strictly speaking
'part of the proof'?'', which produced a useful audit.
\item \emph{Goal-pinning}: dominant. Most interventions were of the form
``yes, but go further'' or ``yes, but match the primer''.
\item \emph{Escalation prompts}: the GPT-5.5 Pro consultations were all
DM-initiated; Claude did not propose them on its own.
\end{itemize}

\section{Observations on the multi-model workflow}

\paragraph{Cost ratio.} Three GPT-5.5 Pro consultations vs.\ ~3000 lines
of Lean. The asymmetry is striking: a few minutes of GPT time produced
the strategic skeleton; many hours of Claude time filled it in. This
matches the pattern in the \texttt{coevolving\_coefficients}
retrospective: GPT was a strategic checkpoint, Claude was the working
partner.

\paragraph{No code from GPT.} Unlike in the grammar-paper collaboration
(where GPT-5.4 Pro produced explicit calculations and ladder-algebra
formulas), in this project GPT-5.5 Pro produced no Lean code. Its
contributions were strategic (which approach), structural (which
hypothesis), and tactical (which Lean idiom). This may reflect the
different shape of the task: formalisation requires less novel
mathematical content and more tactical engineering, and the engineering
is best done in continuous interaction with the Lean compiler --- which
Claude has and GPT does not.

\paragraph{Claude's failure mode: thrashing.} Claude's pattern when
stuck was to produce more variants of the same approach (try
\texttt{linear\_combination}, then \texttt{polyrith}, then
\texttt{nlinarith}, then \texttt{linear\_combination} with hand-crafted
hypotheses) rather than stepping back and asking what the right
\emph{tactic class} for the goal was. The bridge step was eventually
solved by recognising that the goal was a rational identity (so the
right tool is \texttt{field\_simp}) modulo the square-root algebra (so
the right move is to substitute symbols for the square roots). Neither
move requires a flash of mathematical insight; both are standard moves
on a flow chart for ``prove an algebraic identity in Lean''. But Claude
did not make them on its own; the prompt to consult 5.5 Pro was needed.

\paragraph{Claude's overshoot: over-engineering.} Two stages
(\texttt{TailBound} at 532 lines and \texttt{Localisation} at 164 lines)
were over-engineered relative to the actual proof needs. The eventual
proof (per GPT's strategy memo) used a global first-order remainder bound
and never invoked tail bounds at the strength developed. Claude's
defaults pull toward more general infrastructure than is strictly needed.
GPT-5.5 Pro's recommendations consistently pulled in the opposite
direction: ``don't build Watson machinery'', ``skip generic step 10/11'',
``specialise to $n = 0,1,2,3$''. Strategic minimalism was GPT's
contribution; tactical maximalism was Claude's default.

\paragraph{No \texttt{sorry}, ever.} The aks-style discipline held
throughout. At every commit, \texttt{scripts/sorries} reported zero
occurrences. Several intermediate states had stub theorems whose proofs
were sketched in comments but the file was never committed with an
unfilled \texttt{sorry}. This discipline made it possible to interrupt
and resume work without the sense of a growing debt.

\section{Cost and duration}\label{sec:cost}

The full session usage:

\begin{center}
\begin{tabular}{ll}
\hline
Wall-clock duration & 5h 44m 26s \\
API duration & 2h 30m 48s \\
Total cost & \$198.53 \\
Lines added (gross) & 6050 \\
Lines removed (gross) & 1568 \\
Lines retained (net) & $\sim$3279 (final Lean) + $\sim$1200 (markdown, scripts) \\
\hline
\multicolumn{2}{l}{\emph{Per-model usage:}} \\
\hline
Claude (Opus 4.7, 1M context) & \$198.14 \\
\quad input tokens & 6.7k \\
\quad output tokens & 646.2k \\
\quad cache read tokens & 324.2m \\
\quad cache write tokens & 3.2m \\
Haiku 4.5 (auxiliary) & \$0.39 \\
\quad input / output tokens & 72.0k / 5.2k \\
\hline
\end{tabular}
\end{center}

A few features deserve comment.

\paragraph{The wall-clock vs API ratio.} 5h 44m of wall-clock for 2h 30m
of model time means roughly 56\% of the wall-clock was spent waiting on
the human (DM was reading screenshots, drafting messages, deciding what
to do next, or doing other work) and 44\% was spent waiting on the
model. This is a useful sanity check: the bottleneck in this kind of
session is the human, not the model.

\paragraph{Cache read vs input tokens.} 324m cache-read tokens vs 6.7k
fresh input tokens. Almost the entire conversational context was served
from the prompt cache. Without prompt caching this session would have
cost roughly an order of magnitude more.

\paragraph{Lines added vs lines retained.} 6050 lines added, 1568
removed, leaving a final 3279 lines of Lean. So roughly 50\% of what was
written was kept verbatim, 25\% was rewritten, and 25\% was discarded
outright. This is consistent with the impression from the inside: there
were perhaps four or five points where a sub-goal had to be retried or
restructured (the over-engineered \texttt{TailBound} stage, the bridge
step, several \texttt{integral\_add}/\texttt{integral\_sub} pattern
unification failures), and each retry cost a few hundred lines of churn.

\paragraph{Cost per Lean line.} \$198.53 / 3279 retained lines = \$0.06
per kept line of Lean. As a back-of-envelope: at \$50--150/hr for human
formalisation effort and a heuristic of 50--200 lines of Lean per
human-hour for novel content, the equivalent human cost would be
roughly \$800--\$10\,000 for the same output. The advantage compounds
with the human's wall-clock savings: a researcher who would have spent
30 hours on this in 2025 spent under 6 hours of wall-clock attention in
2026, almost all of it directional rather than tactical.

\paragraph{The Haiku rounding error.} The 39 cents of Haiku 4.5 usage was
auxiliary tooling (web fetches, project audits). It is a striking marker
of how cheap small auxiliary models have become, and an argument for
delegating routine searches and audits to them as a matter of course.

\section{Lessons}

\begin{enumerate}
\item \textbf{Borrow infrastructure aggressively.} The
\texttt{aks}-derived audit script, \texttt{lean-search} wrapper, and
\texttt{CLAUDE.md} structure saved at least a day of setup time. A
project of this size cannot afford to grow its own tooling from scratch.

\item \textbf{Strategy consultations are cheap and load-bearing.} Three
short GPT-5.5 Pro consultations changed the trajectory of the entire
project. Without the first, we would have built Watson--Erd\'elyi
infrastructure for weeks. Without the third, we would have stopped at the
$J_n$-form theorem. The \emph{cost} of a strategy consultation is a few
minutes; the \emph{value} is the difference between the project ending
and the project not ending.

\item \textbf{The hypothesis is part of the theorem.} The discriminant
condition $\alpha^2 < 3\lambda\gamma$ is not a regrettable technicality
but a structural feature: without it, the primer's formula is false
(easy counterexample: $\alpha$ very large with $\lambda, \gamma$ small).
The primer treats this as obvious; in Lean it becomes explicit. This is
exactly the bookkeeping benefit DM wanted from formalisation in the
first place.

\item \textbf{Specialise early.} GPT's recommendation to specialise to
$n = 0,1,2,3$ rather than build a generic parity framework was
prescient. Generic frameworks pay off across many uses; for a
single-target proof they are a tax.

\item \textbf{Avoid \texttt{Real.rpow} when integer powers suffice.}
Mathlib's \texttt{Real.rpow} interacts poorly with several common
tactics. For exponents that are known integers, use \texttt{HPow}
directly; the bookkeeping savings compound.

\item \textbf{Notice thrashing early; escalate to a different model.}
DM's intervention to consult 5.5 Pro on the bridge step was the right
move at the right time. Three failed tactical attempts in a row is the
canonical sign. The intervention was not ``DM solved the problem'' but
``DM noticed and changed the configuration''.

\item \textbf{Mathlib gap maps are project-specific.} About 80\% of this
proof was new content built on top of Mathlib's foundations, but the
\emph{shape} of what was missing was specific: pointwise Taylor bounds,
Mill's-ratio family, global Laplace remainder. A different theorem in
SLT would surface a different gap map. The library that this project
contributes (the \texttt{Laplace} library) starts to fill in the
analytic-asymptotics part of that map.

\item \textbf{The point was not the theorem; the point was the loop.}
The covariance formula is not used directly anywhere in Timaeus's
research roadmap. What is used is the loop: ``write LaTeX, formalise the
key claim in Lean, present both to a human reviewer who reads the Lean
statement to verify the LaTeX claim''. This project is one instance of
that loop. Whether the loop scales depends on whether subsequent
projects can be formalised at similar cost; the present project took
roughly a morning of human time, 5h 44m of wall-clock, 2h 30m of API
time, and \$199 of model usage to produce 3279 lines of Lean (see
\S\ref{sec:cost}). If that cost ratio holds, the loop is viable.
\end{enumerate}

\section{Possible next steps}

\paragraph{Tag the primer with a Lean reference.} The primer's
equation~(4.10) ($\Cov_t[x^2,x] = -2\alpha/(\lambda^3 t^2) + o(t^{-2})$)
has been tagged with a \texttt{\textbackslash leanref} macro pointing at
the formalised theorem
\href{https://github.com/timaeus-research/laplace/blob/bfe350c/Laplace/OneD/IntegralRemainder.lean\#L1454}{\texttt{cov\_anharmonic\_asymptotic}},
pinned to commit \texttt{bfe350c}. This is the lightest-weight version
of the loop described in the ``Formalisation First Steps'' note: a
reader of the primer can click through to the Lean statement and verify
in twenty minutes that what is stated in the LaTeX matches what is
proved in Mathlib.

\paragraph{Adopt \texttt{leanblueprint} once a second result is
formalised.} \href{https://github.com/PatrickMassot/leanblueprint}{\texttt{leanblueprint}}
is the de-facto tool for projects of this kind: it is what underpins the
Liquid Tensor Experiment, the Polynomial Freiman--Ruzsa formalisation,
and several other major Mathlib formalisation projects. It provides
LaTeX macros (\texttt{\textbackslash lean\{name\}},
\texttt{\textbackslash leanok}, \texttt{\textbackslash uses\{...\}}) that
mark up paper claims with their corresponding Lean names and dependency
edges, and it generates an interactive HTML dependency graph in which
each node is hyperlinked to the Lean source with build/proof status.
For a single formalised theorem this is overkill --- the
\texttt{\textbackslash leanref} macro above does the job in one line ---
but as soon as a second primer result is formalised (e.g.\
\texttt{lem:laplace\_cov} as the next natural target), the dependency
graph becomes genuinely useful. The setup cost is roughly an afternoon
of work, mostly configuring the blueprint scaffold and porting the
existing \texttt{\textbackslash leanref} call into the
\texttt{\textbackslash lean / \textbackslash leanok / \textbackslash uses}
form.

\paragraph{Multivariate Laplace.} The current formalisation is purely
one-dimensional. The primer's main result
(\texttt{lem:laplace\_cov2} at the multivariate level) is the
$d$-dimensional Wick-contraction formula
$\Cov_t[\phi,\psi] = -\frac{1}{2t^2}\,T_{ijk}[\,\cdots\,] + o(t^{-2})$,
which our \texttt{TailBound.lean} and \texttt{Localisation.lean} were
over-engineered for in anticipation. The 1D case proved that the toolkit
works; the multivariate case will exercise a different part of Mathlib
(\texttt{Matrix} algebra, multivariate Gaussian, contraction with
symmetric tensors).

\paragraph{Backport into Mathlib.} Several pieces of the
\texttt{Laplace} library are general enough to live in Mathlib proper:
the two-sided $|e^{-z} - (1-z)|$ scalar bound
(\texttt{ScalarBound.lean}), the polynomial-times-Gaussian half-line
tail bounds (\texttt{TailBound.lean}), and possibly the rescaling
identity. A PR to Mathlib upstreaming these would shorten future
formalisations and is the kind of contribution that pays off in
goodwill across the Lean community.

\section{Acknowledgements}

\begin{itemize}
\item Geoffrey Irving's \href{https://github.com/girving/aks}{\texttt{aks}}
repository provided the model and several concrete artefacts (audit
script, search wrapper, project layout); see \S\ref{sec:aks}.
\item GPT-5.5 Pro provided three strategic consultations preserved
verbatim at
\href{https://github.com/timaeus-research/laplace/tree/main/gpt_responses}{\texttt{gpt\_responses/}}.
\item The mathematical content tracks the SLT Susceptibility Primer
(Elliott--Murfet, 2026).
\item The formalisation depends throughout on
\href{https://github.com/leanprover-community/mathlib4}{Mathlib}.
\item Edmund Lau and Geoffrey Irving for earlier conversations about
formalisation strategy at Timaeus.
\end{itemize}

\end{document}
